\documentclass[preprint,11pt]{elsarticle}
\usepackage[T1]{fontenc}
\usepackage{lmodern,microtype}
\usepackage{amsmath,amssymb,amsthm,mathtools}
\usepackage{booktabs,tabularx,placeins}
\usepackage[a4paper,margin=26mm]{geometry}
\usepackage[hidelinks]{hyperref}
\newtheorem{theorem}{Theorem}
\newtheorem{proposition}{Proposition}
\newtheorem{corollary}{Corollary}
\newcommand{\E}{\mathbb{E}}
\newcommand{\Prb}{\mathbb{P}}
\newcommand{\R}{\mathbb{R}}
\DeclareMathOperator{\tr}{tr}
\DeclareMathOperator{\Cov}{Cov}
\DeclareMathOperator{\supp}{supp}
\hypersetup{pdftitle={An Exact Pre-Alarm Safety-Accuracy Boundary for Constant-Radius State Saturation}}
\makeatletter
\let\ps@pprintTitle\ps@plain
\makeatother

\begin{document}
\begin{frontmatter}
\title{An Exact Pre-Alarm Safety--Accuracy Boundary for Constant-Radius State Saturation\tnoteref{review}}
\tnotetext[review]{Pre-submission manuscript for internal review.}
\begin{abstract}
State saturation caps the change that an impaired measurement can induce in an estimate. A cap alone leaves open whether a larger radius creates a dangerous estimation displacement before the first alarm. We analyze this question for a fixed finite-duration measurement perturbation, a common trusted conditional mean and a Gaussian mixture Bayes-factor monitor. A joint support condition gives an exact zero-risk boundary within the constant-radius saturation class. Conditional non-concentration resolves equality at the critical radius. A strictly unalarmed support point with two opposite estimate offsets proves positive risk for every larger radius. A first-packet condition tests this support requirement using the innovation covariance and its gain to the state. Nominal conditional covariances then convert the critical radius into a minimum healthy mean-square error. In two frozen five-node Gaussian examples, the minimum nominal increases over the common-gain reference are \(7.070546\%\) and \(6.865134\%\). Deterministic first-packet enclosures and archived paired comparisons quantify positive risk above the boundary. The result concerns the specified saturation class and fixed perturbations.
\end{abstract}
\begin{keyword}
Distributed state estimation \sep finite-horizon safety \sep state saturation \sep Bayes factor \sep multiplicative watermarking
\end{keyword}
\end{frontmatter}

\section{Introduction}

An alarm can arrive after a corrupted measurement has changed a state estimate. This timing matters when the estimate is used during the monitoring window. A bound on the distance between two estimator outputs controls one part of the problem. The data that attain a large displacement must also survive the actual detector up to that time.

We study a prescribed constant-radius saturation around a trusted estimate. The two outputs are coupled through the same physical noise and watermarks. The safety event compares the first crossing of a displacement threshold with the first alarm on the affected received stream. A strict ordering excludes same-round alarms and includes alarms that occur before the first affected packet. The healthy cost is the nominal mean-square error over a fixed horizon.

The analysis establishes when the usual diameter limit is a sharp zero-risk boundary. The converse uses a joint support point under the fixed perturbation law, with both a state displacement in a specified direction and strict survival of the Bayes-factor test. The critical radius yields an exact healthy-error budget test within the saturation class. Two frozen examples support the analysis with certified costs, deterministic first-packet risk bounds and paired counts from previously recorded data. The detector, watermark and estimation procedures are specified inputs to the analysis.

\section{Related work}

Model-consistent faults and attacks have been linked to innovation and residual statistics since the generalized likelihood-ratio work of Willsky and Jones \cite{willsky}. Zhou et al.\ \cite{zhou} use a multiplicative watermark of the form \(\lambda I+\beta\beta^\top\) and analyze residual Kullback--Leibler divergence. Baba and Chattopadhyay \cite{baba} use conditional likelihoods and an onset mixture for quickest detection in correlated distributed tracking. These works motivate the fixed watermark and prefix evidence. Our performance variable additionally depends on two estimation paths and the time of their first separation beyond a prescribed level.

Innovation saturation has established robustness results \cite{he}. Conditional on a trusted history, the radial clipping of an ideal conditional-mean increment has the same one-step form as Ruckdeschel's clipped estimator \cite[Theorem 3.2]{ruck}. This equivalence identifies the saturation map as an existing estimator form. The guaranteed quantities differ. For example, He et al.\ bound errors relative to the physical state under bounded noise and sparse compromised measurements. Li and Mo \cite{limo} compare a secure estimator with nominal or oracle Kalman estimates under sparse corruption. Zhang et al.\ \cite{lee} characterize trade-offs between nominal and adversarial mean-square risks of estimators.

Detector-constrained reachable sets and energy gains already relate detectability to estimation or performance displacement \cite{murguia,gallo}. Joint detection and estimation also has established decision and stopping-risk formulations \cite{sihag,yilmaz}. The boundary studied here concerns the probability of a \emph{strict} first-passage order under one fixed template and one detector. A reachable displacement bound alone does not give its probability under the joint law. The contribution is the support condition that establishes sharpness for this event and connects it to the nominal cost of a fixed-radius class.

\section{System and problem definition}
\label{sec:model}

\subsection{Paired measurements and the reference receiver}

Consider a linear Gaussian state and measurements
\begin{equation}
 X_{k+1}=A X_k+\xi_k,\qquad Y_{j,k}=H_jX_k+\nu_{j,k},\qquad 0\le k<N .
 \label{eq:plant}
\end{equation}
The horizon contains \(N\) updates and starts at index zero. For the frozen examples, the joint physical law is generated by a Gaussian initial state and independent process and measurement noise blocks. The physical variables and the independent legal watermarks form a finite-dimensional Euclidean Borel probability space. Let \(W\) denote the complete watermark sequence on that space. Both paired streams use the same realization of these variables. Conditional on \(W\), only a prescribed perturbation of one received measurement link differs between the streams.

The protected target link uses \(M_k=\lambda I+\beta_k\beta_k^\top\) with \(\lambda>0\), synchronized encoding by \(M_k^{-1}\) and decoding by \(M_k\) \cite{zhou}. In the archived A0 scenario the decoded target measurement is \(Y_s^a=Y_s^0+d_s\) at the first affected index \(s\). The entire sequence of decoded differences is obtained from the single fixed coded template and the realized watermarks. The chosen coded template has finite support and fixed energy. At receiver \(i\) the own measurement and the other neighbor measurement remain shared. Estimates are evaluated on paired trajectories while the physical state equation and the raw monitor are unaffected by the estimator output.

The reference receiver V0 uses the local gain \(K_i\) and the same \(G_i\) for every neighbor in \(\mathcal N_i\). Its update in stream \(v\in\{0,a\}\) is
\begin{equation}
 \widehat X^{\mathrm{V0},v}_{i,k}
   = F_i\widehat X^{\mathrm{V0},v}_{i,k-1}
      +K_iY_{i,k}+G_i\sum_{j\in\mathcal N_i}Y^v_{j\to i,k},
 \qquad F_i=\left(I-K_iH_i-G_i\sum_{j\in\mathcal N_i}H_j\right)A .
 \label{eq:common}
\end{equation}
Initialization in the examples is zero for the estimate at index \(-1\). V1a is the existing structured-compensation reference. It uses \eqref{eq:common} with the posterior mean of the current decoded perturbation subtracted from the target measurement. The posterior is computed from that stream's own legal prefix and the fixed Gaussian working template prior. We use V1a as a strong healthy-accuracy comparator.

\subsection{Trusted center and the saturation class}

Let \(\mathcal T_k\) contain the trusted measurement prefix and the available legal watermarks, and let \(\mathcal F_k\) additionally contain the target prefix under health. Write \(m_k=\E_0[X_k\mid\mathcal T_k]\), \(\mu_k^0=\E_0[X_k\mid\mathcal F_k]\), and let \(\mu_k^a\) be the same nominal conditional-mean map evaluated on the affected input. The nominal map is defined on the received-prefix coordinates. The conditional mean under the true perturbation law need not be this map. Let
\begin{equation}
 g_k=\mu_k^0-m_k,\quad h_k=\mu_k^a-\mu_k^0,\qquad
 \widehat X_k^v(r)=m_k+\Pi_r(\mu_k^v-m_k),
 \label{eq:sat}
\end{equation}
where \(\Pi_r\) is Euclidean projection onto the closed ball of radius \(r\), with \(\Pi_0\equiv0\). This fixed-radius class is the subject of the boundary. It is already implemented in the frozen receiver. For those models \(G_i\) has full row rank, allowing the desired state output to be realized by a target measurement correction in \eqref{eq:common}. No information from the healthy counterfactual is needed by \eqref{eq:sat}.

\subsection{Raw-input alarm and the risk event}

Conditional on a watermark realization, let \(P_0\) be the nominal input law, \(P_a\) the true fixed-template input law, and \(Q\) the normalized working mixture over candidate start times and amplitudes. The Bayes factor \(B_k\) is the likelihood ratio of \(Q\) to \(P_0\) restricted to the raw legal prefix. The likelihoods share the same marginal law for the trusted prefix. In the archive \(Q\) assigns uniform weight over the \(N\) candidate start indices, and each amplitude has the registered zero-mean Gaussian prior. Candidates starting after \(k\) contribute the unit likelihood ratio. The V1a compensation posterior includes the registered no-event prior of \(1/2\). The alarm uses the event-mixture Bayes factor \(B_k\). Both procedures use only the observed prefix and watermarks available at that update.

The first affected-stream alarm occurs when \(B_k^a\ge\eta\); define its time as \(\tau_A\). The first index where \(\|\widehat X_k^a(r)-\widehat X_k^0(r)\|^2\ge D\) is \(\tau_D(r)\), where \(D>0\). An absent crossing or alarm receives the sentinel \(N\). For \(c=\sqrt D\), the strict pre-alarm risk and healthy accuracy loss are
\begin{align}
 R_N(r) &= \Prb\{\tau_D(r)<\tau_A\}
     = \Prb\{\exists k<N:\|\Pi_r(g_k+h_k)-\Pi_r(g_k)\|\ge c,\
                         \max_{j\le k}B_j^a<\eta\},\label{eq:risk}\\
 J_N(r) &= \frac1N\sum_{k=0}^{N-1}\E_0\|X_k-\widehat X_k^0(r)\|^2 .
 \label{eq:mse}
\end{align}
Probability in \eqref{eq:risk} averages the coupled physical and watermark variables with the template fixed. Healthy expectation in \eqref{eq:mse} uses the nominal received stream. The nominal alarm probability is bounded by \(1/\eta\) over the whole horizon under the likelihood-ratio martingale. An alarm before \(s\) therefore still affects the probability in \eqref{eq:risk}.

\section{The zero-risk boundary and its healthy cost}
\label{sec:theory}

Set \(r_\star=c/2\). The two projection images in \eqref{eq:sat} have the same center, giving the standard geometric bound
\begin{equation}
 \|\Pi_r(g_k+h_k)-\Pi_r(g_k)\|\le\min(2r,\|h_k\|).
 \label{eq:diam}
\end{equation}
Consider three conditions on the laws and maps in \eqref{eq:risk}.
\begin{enumerate}
 \item Conditional on \(W\), \(h_k\) is deterministic for each \(k\). On \(\{h_k\ne0\}\), the conditional law of \(g_k\) gives zero probability to \(\operatorname{span}(h_k)\).
 \item The maps \(g_k,h_k,B_k^a\) are continuous on the support of the joint finite input law.
 \item At some index in the monitoring horizon, a single point in that support has strictly unalarmed history. At that point \(g_k\) and \(g_k+h_k\) point in opposite directions and both have length greater than \(r_\star\).
\end{enumerate}
The first condition holds for a conditional Gaussian gap of rank at least two. The last condition is a sufficient joint support test. Its use keeps the alarm and the displacement in the same realization.

\begin{theorem}[Exact boundary within the saturation class]
\label{thm:boundary}
With the shared center in \eqref{eq:sat} and Conditions 1--3,
\begin{equation}
 R_N(r)=0\quad\Longleftrightarrow\quad 0\le r\le r_\star .
 \label{eq:equiv}
\end{equation}
The non-concentration condition in Item 1 suffices for zero risk at equality even without Items 2 and 3.
\end{theorem}

The proof in \ref{app:proof} examines equality in \eqref{eq:diam}. For the converse, both opposite projected radii exceed \(c/2\) at the same strictly unalarmed support point whenever \(r>r_\star\). Continuity then yields a positive-probability neighborhood for each such \(r\). Its probability may depend on \(r\).

\begin{proposition}[First-packet sufficient condition]
\label{prop:prefix}
Suppose the first affected packet is at \(s<N\) and \(\eta>1\). Condition on the full old observations and the current trusted measurements, and denote the nominal innovation covariance of the current target by \(V\succ0\). Let \(L\) be its regression gain to the full state conditional mean. Assume a fixed affine Gaussian nominal model, with legal joint support allowing the centered history and the terminal innovation at a watermark support point. Assume the working mixture has the nominal trusted marginal and unit-covariance shifted Gaussian conditional components after whitening. Then a first-packet displacement \(d\) in the decoded-template support that satisfies
\begin{equation}
 \|Ld\|^2>D,\qquad d^\top V^{-1}d<8\log\eta
 \label{eq:check}
\end{equation}
implies Condition 3. The support requirement is satisfied in the frozen models by noise independent of the watermarks and nondegenerate target innovations.
\end{proposition}

The factor \(8\) follows from evaluating the normalized Gaussian likelihood envelope at half of the first target shift. The assumptions used in this calculation are specified in \ref{app:proof}. The check requires only a nonzero effective state displacement in the tested direction of \(L\). For a fixed first coded template entry \(a\ne0\) and a full-support watermark of the form in Section~\ref{sec:model}, the support contains the ray \(d=\gamma a,\ \gamma\ge\lambda\). Thus one verifiable parameter region is
\begin{equation}
 \lambda^2a^\top V^{-1}a<8\log\eta,\qquad
 0<D< 8\log\eta\,\frac{\|La\|^2}{a^\top V^{-1}a}.
 \label{eq:region}
\end{equation}
These are sufficient inequalities for the registered perturbation law. The abstract converse can also follow from other compatible support points. For instance, a bounded decoded law with all \(\|h_k\|<c\) fails the converse regardless of the detector. This case and the rank-one issue at equality are detailed in \ref{app:proof}.

\begin{corollary}[Minimum nominal cost of zero risk]
\label{cor:cost}
Under Theorem~\ref{thm:boundary}, let \(P_{F,k}\) and \(P_{T,k}\) be the deterministic full and trusted conditional state covariances in the nominal linear Gaussian model. Then
\begin{align}
 J_N(r) &= \frac1N \sum_{k=0}^{N-1}
   \left[\tr P_{F,k}
   +\E \big(\|Z_k\|-r\big)_+^2\right],
       \quad Z_k\sim\mathcal N(0,P_{T,k}-P_{F,k}),\label{eq:cost}\\
 J_{\mathrm{zero}}(D)&=J_N(\sqrt D/2).
 \label{eq:zero-cost}
\end{align}
For \(\epsilon\ge0\), a radius in this class attains zero strict pre-alarm risk and healthy MSE at most \((1+\epsilon)J_{\mathrm{V0}}\) if and only if
\begin{equation}
 J_N(\sqrt D/2)\le (1+\epsilon)J_{\mathrm{V0}}.
 \label{eq:budget}
\end{equation}
\end{corollary}

The cost is continuous, convex and nonincreasing in \(r\). If any gap covariance is nonzero, the cost decreases strictly for every finite radius. If the support condition for the converse is unavailable, a verified diameter condition including equality still gives a sufficient budget test. If equality is unavailable as well, the right endpoint of the strict safety interval may only be an infimum of these certified costs. All statements concern the class \eqref{eq:sat}.

\section{Verification on two frozen models}
\label{sec:evidence}

\subsection{Evaluation protocol and certified operating points}

Both models use a five-node cycle and two-dimensional states. Nodes are indexed from zero. Receiver 1 uses the target link \(0\to1\) and trusted measurements from nodes 1 and 2. The matrices and data are identified in \ref{app:reproduce}. The unstable and stable examples differ in the plant matrix \(A\) and the associated common gains. Both use \(N=32\), the registered ten-packet A0 template starting at \(s=7\), coded energy 10, \(D=0.1\) and \(\eta=20\). The conditional likelihood-ratio calibration gives a nominal event false-alarm bound of \(0.05\).

In each model the state gap covariance is positive definite at every update. The first-packet sufficient ceilings for \(D\) in \eqref{eq:region} are approximately \(0.933871\) and \(0.920396\). The base squared innovation lengths are below \(8\log 20\) in both cases. This verifies the conditions at \(D=0.1\), with \(r_\star=0.1581138830\ldots\). Theorem~\ref{thm:boundary} therefore applies to the 32-update monitoring horizon, as well as the first-packet and first-three-packet prefixes that include \(s\).

\begin{table}[htbp]
\centering
\caption{Nominal mean-square costs and the minimum cost of zero pre-alarm risk. The percentages are relative to V0.}
\label{tab:critical}
\begin{tabular}{lrrr}
\toprule
Plant & V0 & \(J_N(r_\star)\) & Increase (\%) \\
\midrule
Unstable & 0.227711176953 & 0.243811601333 & 7.070546 \\
Stable   & 0.226665590336 & 0.242226485995 & 6.865134 \\
\bottomrule
\end{tabular}
\end{table}

The two values in Table~\ref{tab:critical} use rational conditional covariance recursions and a ball-arithmetic integral. An independent batch covariance calculation and nonnegative tail integral give the same displayed values. In both cases a healthy MSE increase of at most \(2\%\) excludes every zero-risk radius in \eqref{eq:sat}. This conclusion uses the converse in Theorem~\ref{thm:boundary}. With only \eqref{eq:diam}, failure to meet a budget inside the certified safe interval would leave larger radii undecided.

\begin{table}[htbp]
\centering
\small
\caption{Frozen saturation radii and deterministic enclosures of strict first-packet risk, in percent. The two nominal-match radii use a historical point estimate of V1a's healthy cost.}
\label{tab:cert}
\begin{tabular}{lrrrr}
\toprule
Cost point & \multicolumn{2}{c}{Unstable} & \multicolumn{2}{c}{Stable}\\
\cmidrule(lr){2-3}\cmidrule(lr){4-5}
 & Radius & Risk interval & Radius & Risk interval\\
\midrule
\(0.1\%\) & 0.329784 & [0.9062, 7.5011] & 0.326252 & [3.4398, 38.2156]\\
\(0.5\%\) & 0.306555 & [0.6726, 6.7628] & 0.303103 & [2.6762, 35.0549]\\
\(1.0\%\) & 0.284020 & [0.4664, 6.0162] & 0.280671 & [1.9610, 31.8273]\\
\(2.0\%\) & 0.250501 & [0.2164, 4.8585] & 0.247332 & [0.9970, 26.7390]\\
Nominal match & 0.341577 & [1.0289, 7.8618] & 0.342168 & [3.9639, 40.2862]\\
\bottomrule
\end{tabular}
\end{table}

The lower bounds in Table~\ref{tab:cert} include noise regions where a displacement crossing occurs before an alarm. The upper bounds allow every possible displacement crossing, including some that would already have been alarmed. The resulting outer enclosures are computed as described in \ref{app:cert}. The positive lower bounds certify a nontrivial risk at the reported fixed radii.

\subsection{Archived paired evidence and the matching limit}

The held-out archive was generated with independent physical, watermark and candidate onset draws. Conditioning it on the prescribed onset \(7\) and end index \(17\) yields 670 and 681 records. In each record the methods share the input and the original alarm. The first-packet column in Table~\ref{tab:counts} requires the first crossing at \(7\) and a strictly later alarm. The three-packet column accepts a first crossing through \(9\), still strictly before the alarm. The recount applies the same sentinel and tie conventions as \eqref{eq:risk}.

\begin{table}[htbp]
\centering
\caption{Archived event counts on the onset-7 subset. Denominators are 670 for the unstable model and 681 for the stable model.}
\label{tab:counts}
\begin{tabular}{lrrrr}
\toprule
Receiver/operating point & \multicolumn{2}{c}{Unstable} & \multicolumn{2}{c}{Stable}\\
\cmidrule(lr){2-3}\cmidrule(lr){4-5}
 & First & First three & First & First three\\
\midrule
V1a           & 49 & 168 & 263 & 360\\
Saturation \(0.1\%\) & 35 & 123 & 150 & 253\\
Saturation \(0.5\%\) & 30 & 109 & 127 & 221\\
Saturation \(1.0\%\) & 26 & 103 & 114 & 205\\
Saturation \(2.0\%\) & 21 & 80  & 96  & 164\\
\bottomrule
\end{tabular}
\end{table}
\FloatBarrier

At the fixed \(0.1\%\) point the estimated three-packet differences, V1a minus saturation, are \(6.72\) and \(15.71\) percentage points. Bonferroni-controlled paired intervals are \([0.77,12.47]\) and \([8.01,22.99]\) percentage points, under the registered synthetic sampling model. These are comparisons at specified radii. They use discordant pairs rather than an independence assumption between methods.

The exact equal-cost comparison remains unresolved. For V1a the archived healthy conditional scores yield the following cost and corresponding matching-radius intervals, respectively.
\begin{center}
\begin{tabular}{lcc}
\toprule
Plant & Healthy cost & Matching radius \\
\midrule
Unstable & $[0.227493927,\,0.229964382]$ & $[0.284441204,\,0.343889259]$ \\
Stable   & $[0.226356948,\,0.228825678]$ & $[0.282562253,\,0.343587431]$ \\
\bottomrule
\end{tabular}
\end{center}
The held-out outcome files cover the four fixed cost points. Matched-radius evidence comes from the historical point comparisons at the old nominal matches. The repeated Sobol quadrature reproduces those estimates, with scramble variability reported as a numerical diagnostic. The statistical protocol appears in \ref{app:archive}.

\subsection{Scope of the findings}

The certified cost is for a paired estimation displacement before the first alarm. Large errors that both trajectories share are outside this event. The examples share a sensor structure and differ in dynamics, so they establish two parameter checks of the theorem. They do not establish a network-wide theorem over arbitrary disturbed links, watermarks or disturbances. In particular, changing the trusted information, recoupling the detector to the estimator, or changing the prefix Bayes-factor law can change the support test and the healthy cost. Both receiver classes here are pre-existing.


\section{Conclusion}

A common trusted center and conditional non-concentration give a zero pre-alarm risk guarantee at the diameter threshold. A compatible unalarmed support point makes that threshold exact. In the two examples the healthy MSE budget of \(2\%\) is below the minimum required for zero risk in the fixed-radius class. The archived comparisons resolve some fixed operating points. Performance ordering at exactly equal nominal cost remains open for these data. The boundary depends on the stated information, support and stopping assumptions, and finite horizons with a changed monitoring rule require separate verification.

\section*{Data and code availability}

The model, synthetic records, receivers and input manifests are available in the
\href{https://github.com/center-666/paperlast/tree/b1509be0265ecd0512befe999847fa5908087872}{frozen research archive}.
The accompanying pre-submission audit provides a read-only batch-covariance and archived-event check. The arithmetic and evidence scope are described in \ref{app:cert} and \ref{app:archive}. Code locations are listed in \ref{app:reproduce}.

\appendix
\section{Proof details and scope of the first-packet test}
\label{app:proof}

\subsection{Critical equality and the open-neighborhood argument}

The two outputs in \eqref{eq:sat} differ only by their projected increments. The distance between them is at most \(2r\) and also at most \(\|h_k\|\) by nonexpansiveness. At \(r=c/2\), attaining \(c\) forces both projected increments to have length \(r\) and to be antiparallel. If \(h_k=0\) this is impossible. Otherwise projection preserves the direction of each nonzero increment, so \(g_k\) belongs to the line spanned by \(h_k\). Condition on \(W\), apply the non-concentration requirement, and integrate. Each such event has zero probability. Taking a union over the finite horizon proves the claim at the critical point.

For \(r>c/2\), let the two opposite lengths at the support point of Condition 3 be \(a_0,a_1>c/2\). The projected distance there is \(\min(a_0,r)+\min(a_1,r)>c\). The strict threshold and strict unalarmed prefix define a relatively open subset of the support containing this point, hence a set of positive probability. This establishes Theorem~\ref{thm:boundary}. Bounded watermarks can keep \(\|h_k\|<c\) throughout a finite horizon, in which case nonexpansiveness gives zero risk for all \(r\). Conversely, if \(g_k=Ue_1\) and \(h_k=be_1\) with \(b>c\), then \(U\in[r_\star-b,-r_\star]\) attains the diameter at \(r_\star\). A positive mass of unalarmed outcomes on this interval invalidates equality safety. These cases identify distinct failures of the support and non-concentration conditions.

\subsection{Conditional law, joint support and the likelihood envelope}

Let \(\mathcal H_s\) comprise the old full received prefix and current trusted measurements. The old target prefix is healthy under the true fixed-template law. Conditional on \(W\), subtract the nominal affine regression of the healthy current target on \(\mathcal H_s\). The resulting innovation has conditional mean zero and covariance \(V\). At the mean of the history the trusted and the full-old-plus-current-trusted state means coincide. At a joint support point with healthy terminal innovation \(-d/2\), the two full-mean increments relative to that center are \(-Ld/2,Ld/2\). Independence of the Gaussian physical vectors and \(W\), the mean's presence in the support of a Gaussian history, and \(V\succ0\) justify the point in the frozen model. If a displacement is specified by a limit point of the watermark image, the strict inequalities in \eqref{eq:check} allow a nearby actual support value by continuity.

For a prefix, whiten the nominal target observations given the trusted observations, in time-ordered blocks. Under each fixed working template parameter the conditional whitened innovation vector has a shifted unit-covariance Gaussian density. Under their normalized mixture its density is bounded by the height of the standard Gaussian at zero. The ratio to the nominal standard Gaussian density is bounded by \(\exp(\|z\|^2/2)\). In particular, the Bayes factor at the centered unaffected prefixes is at most one. At the chosen first-packet support point, the affected vector has earlier blocks zero and final squared length \(d^\top V^{-1}d/4\). Thus \eqref{eq:check} implies a strictly unalarmed prefix and both lengths exceeding \(c/2\). Nonzero nominal means can be subtracted. Independence here provides joint support; the abstract theorem only uses the actual joint support. The derivation also uses common trusted marginals under the working mixture and absence of dependence of the alarm on repaired estimates.

\subsection{Nominal MSE identity}

Given \(\mathcal F_k\), the conditional expectation of \(X_k-\mu_k^0\) is zero. Expanding the squared error in \eqref{eq:sat} gives
\[
 \E_0\|X_k-\widehat X_k^0(r)\|^2
 =\tr P_{F,k}+\E_0\|\mu_k^0-\widehat X_k^0(r)\|^2 .
\]
The second norm is the distance from the full mean to the translated ball, namely \((\|g_k\|-r)_+\). By the tower property the conditional mean gap has zero mean and covariance \(P_{T,k}-P_{F,k}\). The gap is Gaussian under the stated deterministic affine information model. For any fixed \(x\ge0\), \((x-r)_+^2\) is convex in \(r\) and differentiable with derivative \(-2(x-r)_+\). Dominated convergence gives the derivative of \eqref{eq:cost}. If a gap covariance is nonzero, its Gaussian norm has unbounded support and the expected derivative is strictly negative for finite \(r\). Minimization over \([0,r_\star]\) now gives \eqref{eq:zero-cost} and \eqref{eq:budget}.

\section{Deterministic numerical enclosures}
\label{app:cert}

\subsection{Healthy costs}

The nominal covariances use decimal model entries interpreted exactly as rationals. The trace of the common-gain reference covariance serves as the V0 denominator. For the two-dimensional gap, let \(\lambda_-\le\lambda_+\) be the covariance eigenvalues and \(v(\theta)=\lambda_-+(\lambda_+-\lambda_-)\sin^2\theta\). Conditional on direction, the standard Gaussian radius has a Rayleigh distribution. Its truncated second moment is
\[
 f(v,r)=2v e^{-r^2/(2v)}-r\sqrt{2\pi v}\,
             \operatorname{erfc}\!\left(r/\sqrt{2v}\right).
\]
One obtains this expression by integrating \(2u\exp(- (r+u)^2/(2v))\) for \(u\ge0\). The expectation is \(2/\pi\) times the integral of \(f(v(\theta),r)\) on \([0,\pi/2]\). The rational recursions and certified angle integrals give the respective critical MSE enclosures
\begin{align*}
 \text{Unstable:}\quad &[0.24381160133333157,\,0.24381160133333163],\\
 \text{Stable:}\quad &[0.24222648599508800,\,0.24222648599508806].
\end{align*}
The independent floating-point check forms the entire finite joint state covariance, applies a Schur complement to each measurement prefix, and integrates the nonnegative radial tail before averaging the directions. This check uses a different covariance and integration calculation from the certificate.

\subsection{First-packet risk lower and upper bounds}

With \(V=C_V C_V^\top\), write \(g_s=g_H+LC_V U\) and \(\mu=C_V^{-1}d_s\), where \(U\) is a standard Gaussian independent of the old and trusted data conditional on the legal watermark. The residual of this conditional split, \(g_H\), has covariance \(P_{T,s}-P_{H,s}\), with \(P_{H,s}\) the state covariance given \(\mathcal H_s\). Earlier false alarms have probability at most \(1/\eta\) even after conditioning on \(U\). For the current alarm, write the Bayes factor as a marginal history likelihood ratio \(B_H\) times the current conditional ratio. The density-peak inequality gives \(B_s\le B_H e^{\|U+\mu\|^2/2}\), with nominal \(\E_0[B_H\mid W]=1\). The history is actually unaffected by the template at the first packet. A union bound therefore yields
\[
 \Prb(\tau_A\le s\mid W,U)\le \eta^{-1}\big(1+e^{\|U+\mu\|^2/2}\big).
\]
This bound retains the possibility of earlier and same-round alarms. It does not use independence of the danger and alarm events.

For a box in two watermark and two terminal-innovation coordinates, let \(a\) be a certified radius such that all points in the box are dangerous whenever \(\|g_H\|\le a\). Set \(\sigma_H^2=\lambda_{\max}(P_{T,s}-P_{H,s})\). The conditional risk on this box is at least
\[
 \left[1-e^{-a^2/(2\sigma_H^2)}
     -\eta^{-1}\big(1+e^{\|U+\mu\|^2/2}\big)\right]_+.
\]
The Rayleigh bound follows by representing \(g_H\) as a linear image of a two-dimensional standard Gaussian. Possible boxes for the upper bound use instead the two Gaussian coordinates of the entire healthy gap \(g_s\). Any point with an output displacement beyond \(c\) is included. The radius and displacement tests include
\[
 \|\Pi_r(u)-\Pi_r(v)\|
 \le \frac{2r}{\max(r,\|u\|)+\max(r,\|v\|)}\,\|u-v\|.
\]
When both points are outside, squaring reduces the inequality to
\((\|u\|-\|v\|)^2(1+\cos\theta)\ge0\). When one point is inside, the squared inequality is affine in \(\cos\theta\) and holds at its endpoints \(1\) and \(-1\) since the inside length is at most \(r\) and the outside length at least \(r\). Inside both, the inequality is equality.

Each truncated integration uses a Cartesian mesh of step \(1/16\). The half-grid for the first watermark coordinate uses sign symmetry. The watermark cutoffs are \(2.75\) and the two standard Gaussian cutoffs \(3.5\), giving \(48{,}570{,}368\) candidate cell pairs. Cell masses are enclosed by Gaussian CDF differences. Their omitted joint mass is at most \(0.001131346363\), charged completely to the upper bound. State projection distances have a \(10^{-8}\) binary64 midpoint margin after interval propagation of the box coordinates. The code checks coordinate and radius ranges. The nonnegative summation margin is \(10^{-7}\). These margins enclose rounding under correctly rounded IEEE~754 binary64 operations, including the square root, and the stated ball-arithmetic library. The calculations produce outer enclosures, with analytic relaxations for alarms and for partially classified boxes.

\section{Statistical evidence and cost matching}
\label{app:archive}

All methods share archived inputs in the held-out comparisons. For a contrast, let \(p_+\) and \(p_-\) be the two marginal probabilities of a discordant outcome in one record. The risk contrast is \(p_+-p_-\). Single-method probabilities and these discordant probabilities have binomial counts across records. The 52 reported binomial intervals share a Bonferroni error budget of \(0.025\); differences use the corresponding outward endpoint subtraction. The full archive includes an unobserved-event sentinel and the strict comparison excludes ties.

The healthy V1a score is the time average of \(\|\widehat X_k^{\rm V1a}-\mu_k^0\|^2\), plus the deterministic full-information covariance trace average for the MSE. For each model there are 6144 saved scores from two independent registered families. The lower mean interval uses a mixture of nonnegative products for a nonnegative mean. The upper interval uses bounded products after clipping the scores at \(2.6\), then adds a deterministic upper bound for the excess moment. The latter follows from a rank-two Gaussian quadratic envelope with at most \(32(10+1)\) terms. Each MSE endpoint has error probability at most \(0.00625\). The complete health family therefore has budget \(0.025\), and together with the risk family the total error budget is \(0.05\) under the registered sampling law. Radius-endpoint inclusions use the monotonicity of \(J_N\) and certified sign checks. The intervals describe statistics of this recorded score implementation; the critical zero-risk costs themselves were obtained deterministically.

The historical deterministic-input Sobol rerun uses four fixed scrambles with nested sizes 512 and 1024. Its three-packet point estimates, V1a and nominal-match saturation respectively, are \(24.19\%,18.63\%\) and \(54.39\%,40.82\%\). Variation across scrambles is an integration diagnostic. The archived count intervals and the deterministic first-packet enclosures have separate meanings.

\section{Frozen model and reproduction index}
\label{app:reproduce}

The model in Section~\ref{sec:evidence} uses \(X_0\sim\mathcal N(0,I_2)\), process covariance \(I_2\), and independent measurement covariances \(0.5I_2\). The two plant matrices are
\[
 A_{\rm u}=\begin{bmatrix}1.08&0.05\\0.01&1.05\end{bmatrix},
 \qquad
 A_{\rm s}=\begin{bmatrix}0.80&0.05\\0.01&0.90\end{bmatrix}.
\]
The retained model file lists all observations, neighbor sets and common gains. The watermark uses \(\lambda=0.5\) and independent \(\beta_k\sim\mathcal N(0,0.5I_2)\). The supplied working prior and sequence of starts are defined in the original registration. Case results use the original energy-normalized observation template in the archive.

\begin{itemize}
 \item \href{https://github.com/center-666/paperlast/tree/b1509be0265ecd0512befe999847fa5908087872/common_g_reference_revision_20261115/inputs}{Model and common gains}.
 \item \href{https://github.com/center-666/paperlast/tree/b1509be0265ecd0512befe999847fa5908087872/JFI_secure_estimation_accuracy_safety_frontier}{Receiver and archived counts}.
 \item \href{https://github.com/center-666/paperlast/tree/b1509be0265ecd0512befe999847fa5908087872/JFI_A0_safety_accuracy_theorem_review}{Health, first-packet certificates and archived confidence coverage}.
 \item \href{https://github.com/center-666/paperlast/tree/b1509be0265ecd0512befe999847fa5908087872/JFI_A0_direct_joint_risk_feasibility}{Fixed quadrature and paired-output recurrence audit}.
\end{itemize}
File digests and the independent presubmission check accompany the manuscript. All evaluations are offline and use existing records or the frozen analytic integrals.

\begin{thebibliography}{99}
\raggedright

\bibitem{willsky}
A. S. Willsky and H. L. Jones,
A generalized likelihood ratio approach to the detection and estimation of jumps in linear systems,
\emph{IEEE Transactions on Automatic Control} 21 (1) (1976) 108--112.
\href{https://doi.org/10.1109/TAC.1976.1101146}{doi:10.1109/TAC.1976.1101146}.

\bibitem{zhou}
J. Zhou, W. Yang, W. Ding, W. X. Zheng and Y. Xu,
Watermarking-based protection strategy against stealthy integrity attack on distributed state estimation,
\emph{IEEE Transactions on Automatic Control} 68 (1) (2023) 628--635.
\href{https://doi.org/10.1109/TAC.2022.3171422}{doi:10.1109/TAC.2022.3171422}.

\bibitem{baba}
S. A. Baba and A. Chattopadhyay,
Quickest detection of false data injection attack in distributed process tracking,
\emph{Signal Processing} 239 (2026) 110294.
\href{https://doi.org/10.1016/j.sigpro.2025.110294}{doi:10.1016/j.sigpro.2025.110294}.
Numbering in the audit follows \href{https://arxiv.org/abs/2402.09743}{arXiv:2402.09743}.

\bibitem{he}
X. He, X. Ren, H. Sandberg and K. H. Johansson,
How to secure distributed filters under sensor attacks,
\emph{IEEE Transactions on Automatic Control} 67 (6) (2022) 2843--2856.
\href{https://doi.org/10.1109/TAC.2021.3092603}{doi:10.1109/TAC.2021.3092603}.

\bibitem{ruck}
P. Ruckdeschel,
Optimally (distributional-)robust Kalman filtering,
\href{https://arxiv.org/abs/1004.3393}{arXiv:1004.3393}, 2010.

\bibitem{limo}
Z. Li and Y. Mo,
Efficient secure state estimation against sparse integrity attack for system with non-derogatory dynamics,
\href{https://arxiv.org/abs/2106.03066}{arXiv:2106.03066}, 2021.
Journal version titled \emph{Efficient secure state estimation against sparse integrity attack for regular linear system},
\emph{International Journal of Robust and Nonlinear Control} 33 (1) (2023) 209--236.
\href{https://doi.org/10.1002/rnc.6264}{doi:10.1002/rnc.6264}.

\bibitem{lee}
T. T. C. K. Zhang, B. D. Lee, H. Hassani and N. Matni,
Adversarial tradeoffs in robust state estimation,
\emph{Proceedings of the American Control Conference}, 2023, pp.\ 4083--4089.
\href{https://arxiv.org/abs/2111.08864}{arXiv:2111.08864}.

\bibitem{murguia}
C. Murguia and J. Ruths,
On reachable sets of hidden CPS sensor attacks,
\emph{Proceedings of the American Control Conference}, 2018, pp.\ 178--184.
\href{https://arxiv.org/abs/1710.06967}{arXiv:1710.06967}.

\bibitem{gallo}
A. J. Gallo, S. C. Anand, A. M. H. Teixeira and R. M. G. Ferrari,
Switching multiplicative watermark design against covert attacks,
\emph{Automatica} 177 (2025) 112301.
\href{https://doi.org/10.1016/j.automatica.2025.112301}{doi:10.1016/j.automatica.2025.112301}.
\href{https://arxiv.org/abs/2502.18948v1}{Author preprint}.

\bibitem{sihag}
S. Sihag and A. Tajer,
Secure estimation under causative attacks,
\emph{IEEE Transactions on Information Theory} 66 (8) (2020) 5145--5166.
\href{https://doi.org/10.1109/TIT.2020.2985956}{doi:10.1109/TIT.2020.2985956}.

\bibitem{yilmaz}
Y. Yilmaz, G. V. Moustakides and X. Wang,
Sequential joint detection and estimation,
\emph{Theory of Probability and Its Applications} 59 (3) (2015) 452--465.
\href{https://doi.org/10.1137/S0040585X97T987235}{doi:10.1137/S0040585X97T987235}.
\href{https://arxiv.org/abs/1302.6058}{Author preprint}.

\end{thebibliography}

\end{document}
